\subsubsection{Algoritmo Hungaro (Assignment Problem)}
\begin{lstlisting}[style=mystyle, language=C++]
// TC: O(N^2 * M), si N = M => O(N^3)
// usa: matching bipartito con pesos / problema de asignacion (min cost)
// 1-indexed: n filas (trabajadores), m columnas (tareas), requiere n <= m
// Para MAX cost: pasar -cost[i][j] y el costo maximo es -res.first
#include <bits/stdc++.h>
using namespace std;

const long long INF = 1e18;

// cost[1..n][1..m] es el costo de asignar trabajador i a tarea j
// Retorna {costo_minimo, asignacion}: ans[i] = tarea asignada al trabajador i
pair<long long, vector<int>> hungarian(int n, int m, const vector<vector<long long>>& cost){
  vector<long long> u(n + 1, 0), v(m + 1, 0);
  vector<int> p(m + 1, 0), way(m + 1, 0);

  for(int i = 1; i <= n; i++){
    p[0] = i;
    int j0 = 0;
    vector<long long> minv(m + 1, INF);
    vector<bool> used(m + 1, false);

    do {
      used[j0] = true;
      int i0 = p[j0], j1 = 0;
      long long delta = INF;

      for(int j = 1; j <= m; j++){
        if(!used[j]){
          long long cur = cost[i0][j] - u[i0] - v[j];
          if(cur < minv[j]){
            minv[j] = cur;
            way[j] = j0;
          }
          if(minv[j] < delta){
            delta = minv[j];
            j1 = j;
          }
        }
      }

      for(int j = 0; j <= m; j++){
        if(used[j]){
          u[p[j]] += delta;
          v[j] -= delta;
        } else {
          minv[j] -= delta;
        }
      }
      j0 = j1;
    } while(p[j0] != 0);

    do {
      int j1 = way[j0];
      p[j0] = p[j1];
      j0 = j1;
    } while(j0);
  }

  long long min_cost = -v[0];
  vector<int> ans(n + 1, 0);
  for(int j = 1; j <= m; j++){
    if(p[j] > 0) ans[p[j]] = j; // trabajador p[j] asignado a tarea j
  }
  return {min_cost, ans};
}

int main(){
  int n = 3, m = 3;
  // Matriz de costos 1-indexed de tamano (n+1) x (m+1)
  vector<vector<long long>> cost = {
    {0, 0, 0, 0},
    {0, 4, 1, 3},
    {0, 2, 0, 5},
    {0, 3, 2, 2}
  };

  auto [min_cost, match] = hungarian(n, m, cost);
  cout << "Costo minimo = " << min_cost << "\n"; // 5
  for(int i = 1; i <= n; i++){
    cout << "Trabajador " << i << " -> Tarea " << match[i] << "\n";
  }
  return 0;
}
\end{lstlisting}
